% BHU PYQ -- Professional Exam Template
\documentclass[11pt,a4paper]{article}

\usepackage[a4paper,top=2.5cm,bottom=2.5cm,left=2.8cm,right=2.8cm,headheight=14pt]{geometry}
\usepackage{amsmath,amssymb}
\usepackage[T1]{fontenc}
\usepackage[utf8]{inputenc}
\usepackage{lmodern}
\usepackage{microtype}
\usepackage{fancyhdr}
\usepackage{enumitem}
\usepackage{listings}
\usepackage{hyperref}
\usepackage{lastpage}
\usepackage{parskip}

\hypersetup{colorlinks=true,urlcolor=black,linkcolor=black,
  pdftitle={CS-104b: Computer Org \& Arch -- BHU PYQ}}

\pagestyle{fancy}
\fancyhf{}
\renewcommand{\headrulewidth}{0.4pt}
\renewcommand{\footrulewidth}{0.4pt}
\fancyhead[L]{\small\textit{Banaras Hindu University}}
\fancyhead[R]{\small\textit{CS-104b: Computer Org \& Arch}}
\fancyfoot[C]{\small Page \thepage\ of \pageref{LastPage}}

\lstset{
  basicstyle=\small\ttfamily,
  frame=single,
  breaklines=true,
  columns=flexible,
  keepspaces=true
}

\newlist{parts}{enumerate}{2}
\setlist[parts,1]{label=(\alph*),leftmargin=*,itemsep=4pt,topsep=3pt}
\setlist[parts,2]{label=(\roman*),leftmargin=*,itemsep=2pt,topsep=2pt}
\newcommand{\pts}[1]{\hfill\textit{[#1]}}

\begin{document}

\begin{center}
  {\large\bfseries BANARAS HINDU UNIVERSITY}\\[2pt]
  {\normalsize B.Sc. (Hons.) Semester IV Examination, 2023-24}\\[6pt]
  \rule{0.8\linewidth}{0.5pt}\\[6pt]
  {\large\bfseries Paper: CS-104b --- Computer Org \& Arch}\\[4pt]
  {\normalsize Time: Three hours \hfill Maximum Marks: 70}
\end{center}

\rule{\linewidth}{0.4pt}

\smallskip
\noindent\textit{Instructions: Answer five questions in all including Question No. 1, which is compulsory. Answers should be brief and to the point and may be supplemented with neat sketches, wherever required.}

\bigskip

\medskip\noindent\textbf{Question 1.}
\begin{parts}
  \item A cache memory has 8 cache lines (0-7). The memory block requests are in the order: 4, 3, 25, 8, 19, 6, 25, 8, 16, 35, 45, 22, 8, 3, 16, 25, 7.
  \begin{parts}
    \item Which cache line will have memory block 7 if fully associative mapping with LRU policy is used? Also, calculate the hit ratio and miss ratio.
    \item Which cache line will have memory block 7 if direct mapping is used? Also, calculate the hit ratio and miss ratio.
  \end{parts} \pts{7}
  \item What are the basic differences among a branch instruction, a call subroutine instruction, and a program interrupt? \pts{7}
\end{parts}

\medskip\noindent\textbf{Question 2.}
\begin{parts}
  \item Design a 4-bit Binary Adder-Subtractor and a Binary Incrementor. Explain their working. \pts{6}
  \item Why is register transfer language preferred for describing the internal organization of digital computers? \pts{4}
  \item What is an instruction cycle, and what are the phases of the instruction cycle? \pts{4}
\end{parts}

\medskip\noindent\textbf{Question 3.}
\begin{parts}
  \item Design a bus system using three-state buffers. The number of registers is 6, and the size of each register is 4 bits. \pts{6}
  \item What is the difference between isolated I/O and memory-mapped I/O? \pts{4}
  \item Write the differences between Static RAM and Dynamic RAM. \pts{4}
\end{parts}

\medskip\noindent\textbf{Question 4.} A computer uses a memory unit with 1 Giga words of 16 bits each. The computer supports three addressing modes (register mode, register indirect mode, direct mode) and 35 unique micro-operations. The instruction has four parts: an addressing mode part, an operation code, a register code part to specify one of 180 registers, and an address part.
\begin{parts}
  \item
  \begin{parts}
    \item Find the minimum number of words required for storing one instruction.
    \item Draw the instruction word format and indicate the number of bits in each part.
  \end{parts} \pts{6}
  \item Explain any four addressing modes with suitable examples. \pts{4}
  \item Explain any four applications of the stack in computer organization. \pts{4}
\end{parts}

\medskip\noindent\textbf{Question 5.}
\begin{parts}
  \item What is virtual memory? Explain Virtual address Mapping using Pages with necessary examples. \pts{7}
  \item What is the locality of reference in cache memory? Explain its major types and how they are achieved. \pts{4}
  \item Explain word addressable and byte addressable memory. \pts{3}
\end{parts}

\medskip\noindent\textbf{Question 6.}
\begin{parts}
  \item We have to provide a memory capacity of 4096 bytes using $128 \times 8$ RAM chips.
  \begin{parts}
    \item How many chips are required?
    \item How many lines of the address bus must be used to access 4096 bytes of memory?
    \item How many of lines in (ii) will be common to all chips?
    \item How many lines must be decoded for chip select? Specify the size of decoders.
  \end{parts} \pts{5}
  \item What is an interrupt? Explain its major types. \pts{5}
  \item List out the importance of the I/O interface. \pts{4}
\end{parts}

\medskip\noindent\textbf{Question 7.}
\begin{parts}
  \item Explain the working of DMA (Direct Memory Access) with the block diagram. \pts{8}
  \item What is the disadvantage of the strobe method? Explain how the handshake method solves the problem. \pts{6}
\end{parts}

\vfill
\rule{\linewidth}{0.4pt}
\begin{center}\small
\href{https://bhu-exam.vercel.app/}{Click here to visit our website}\\[4pt]\href{https://me-aryan.vercel.app/}{Click here to visit our contributor}\\[4pt]\href{https://quashan-me.vercel.app/}{Click here to visit our contributor}
\end{center}
\end{document}
